\documentclass[11pt]{article}
\usepackage[margin=1.15in]{geometry}
\usepackage{amsmath,amssymb,amsthm}
\usepackage{hyperref}
\sloppy
\newtheorem{theorem}{Theorem}
\newtheorem{proposition}[theorem]{Proposition}
\newtheorem{lemma}[theorem]{Lemma}
\newtheorem{corollary}[theorem]{Corollary}
\newtheorem{conjecture}[theorem]{Conjecture}
\theoremstyle{remark}\newtheorem{remark}[theorem]{Remark}
\DeclareMathOperator{\perm}{per}
\newcommand{\E}{\mathbb E}

\title{The permanent of the GCD matrix}
\author{Vico Bonfioli\thanks{Data and scripts: \texttt{code/} in the
\texttt{rama-notes} repository; an independent Lean/Mathlib formalization of
Theorem~\ref{thm:cong}'s engine is in \texttt{RamaLean/Paper3Permanent.lean}. Correspondence: \texttt{vicobonfioli@gmail.com}.}}
\date{July 2026}

\begin{document}
\maketitle

\begin{abstract}
Let $a(n)=\perm[\gcd(i,j)]_{1\le i,j\le n}$. Whereas the determinant
$\det[\gcd(i,j)]=\prod_{k\le n}\varphi(k)$ has been known since Smith (1876),
the permanent---sequence \texttt{A085244} in the OEIS---appears not to have been
studied beyond its definition and values. We compute $a(n)$ for
$n\le 30$, prove a family of congruences---$p\mid a(n)$ once there are at least
$p$ integers $\le n$ all of whose prime factors are $\equiv 1\pmod p$ (giving
$2\mid a(n)$ for $n\ge 3$, sharpened to $4\mid a(n)$ for $n\ge4$, and
$3\mid a(n)$ for $n\ge 13$), and show $v_2(a(n))\to\infty$ via a
$(\mathbb Z/2)^k$ group action on the ``$\equiv1\bmod 2^k$'' rows---and prove
$(a(n)/n!)^{1/n}\to\infty$: writing $(a(n)/n!)^{1/n}=(\log n)^{\theta+o(1)}$, the
exponent obeys $2\ln\varphi-\varphi^{-2}\le\theta\le\log2$ (numerically
$\theta\approx0.68$; whether $\theta=\log2$ or a strictly smaller constant is
open). The upper bound $\theta\le\log2$ is
elementary; the lower bound is a doubly stochastic construction from
prime-divisibility blocks (via Van der Waerden--Falikman--Egorychev) whose
per-prime gains $g_p$ satisfy $p\,g_p\to 2\ln\varphi-\varphi^{-2}$ (with
$\varphi$ the golden ratio), so
$\sum_pg_p=\infty$ by Mertens \cite{HW}---giving $c_\infty=\infty$ unconditionally.
Finally, we explain the $2$-adic \emph{deficit} $v_2(n!)-v_2(a(n))$ at $n=2^k+1$ through a
mechanism rather than opaque cancellation: two exact valuation theorems for odd permanents (of size
$2^a$, resp.\ $2^a+1$, driven by the first, resp.\ second, permanental symmetric function) combine to
give $D_{N_0}(2^k+1)=2k-4$ --- for the \emph{weight-zero grade} $N_0$ of $a$, not for $a$ itself ---
whenever a computable, permanent-free ``achiever parity'' is odd, which occurs for $11$ of the $20$
values $k\le23$. So that grade's deficit is $\Theta(\log n)$ conditional on the parity's infinitude,
one open step; transferring the conclusion to $a$ itself needs $v_2(a)=v_2(N_0)$, which we verify only
at $k\le5$ and where the grade margins are small, so we do not claim it in general. Reducing that parity through a chain of verified identities
(machine-checked at five links) yields the following description: the boundary contribution of each
prime $p\equiv3\ (4)$ is exactly the \emph{binary expansion of $1/p$}, and the resulting words'
statistics form a class-number ladder of \emph{depth exactly three}---density, digit-pair, and
digit-triple counts are exact formulas in $h(-p)$ and $h(-8p)$ (e.g.\ the number of $11$-pairs per
period is $(p-4h(-p)-3)/8$), while digit-quadruple and finer statistics provably admit no
class-number formula, the ladder terminating because no real primitive Dirichlet character of
conductor $16$ exists. The $2$-adic deficit of the gcd permanent is thus completely stratified: a
finite class-number core over $\mathbb Q(\sqrt{-p})$ and $\mathbb Q(\sqrt{-2p})$, and a
character-sum (Pólya--Vinogradov) tail.
\end{abstract}

\section{The sequence}
$a(n)=\sum_{\sigma\in S_n}\prod_{i=1}^n\gcd(i,\sigma(i))$ begins
$1,3,14,112,872,14372,154480,\dots$ (30 terms computed exactly, by a Gray-code
Ryser evaluation \cite{Ryser} and, for $n>27$, a modular Ryser with CRT). It is sequence
\texttt{A085244} in the OEIS (Dekel, 2003, with a $b$-file to $n=38$); its
arithmetic---the subject of this note---appears not to have been studied, and the
derived $2$-adic valuation, deficit, and achiever-parity sequences are new.
Computing permanents is $\#P$-hard (Valiant), and---unlike the
determinant---there is no product formula.

\section{Congruences}
For an arithmetic function the determinant $[\,f(\gcd(i,j))\,]$ has Smith's
product form; the permanent has no such structure, but retains congruences.

\begin{proposition}\label{prop:par}
$a(n)\equiv\prod_{k\le n}\varphi(k)\pmod 2$; in particular $2\mid a(n)$ for
$n\ge 3$.
\end{proposition}
\begin{proof}
Over $\mathbb F_2$ permanent and determinant agree, so $a(n)\equiv
\det[\gcd(i,j)]=\prod_{k\le n}\varphi(k)$; the factor $\varphi(3)=2$ makes this
even for $n\ge 3$.
\end{proof}

\begin{theorem}\label{thm:cong}
Let $p$ be prime and $A_p(n)=\#\{j\le n:\text{every prime factor of }j\equiv
1\ (\mathrm{mod}\ p)\}$. If $A_p(n)\ge p$ then $p\mid a(n)$. Consequently
$2\mid a(n)$ for $n\ge 3$ and $\mathbf{3\mid a(n)}$ \textbf{for all} $n\ge 13$
(the third such $j$ is $13$); more generally $5\mid a(n)$ for $n\ge 61$, etc.
\end{theorem}
\begin{proof}
\emph{Lemma A.} If a matrix $B$ over $\mathbb Z$ has $p$ columns pairwise
congruent mod $p$, then $\perm(B)\equiv 0\pmod p$. Indeed, let $T$ index $p$
such columns (common value $v\bmod p$); group the permutations by the pair
$(R,\rho)$ with $R=\sigma^{-1}(T)$ and $\rho=\sigma|_{[n]\setminus R}$. Exactly
$p!$ permutations give each $(R,\rho)$, and for all of them
$\prod_i B_{i,\sigma(i)}\equiv(\prod_{i\in R}v_i)\prod_{i\notin R}B_{i,\rho(i)}
\pmod p$; summing, $\perm(B)\equiv p!\sum(\cdots)\equiv 0$.

\emph{Lemma B.} As $i$ runs over $[1,n]$, $\gcd(i,j)$ takes every divisor of
$j$; so the $j$-th column of $[\gcd(i,j)]$ is $\equiv(1,\dots,1)\pmod p$ iff
every divisor---equivalently every prime factor---of $j$ is $\equiv1\pmod p$.

If $A_p(n)\ge p$ there are $\ge p$ such columns, all $\equiv(1,\dots,1)\bmod p$;
apply Lemma A. The counts follow from the $p=2$ ($j$ odd) and $p=3$
($j\in\{1,7,13,\dots\}$) cases.
\end{proof}
Theorem~\ref{thm:cong} is formalized \emph{end to end} in Lean 4/Mathlib with
only the three standard axioms (no \texttt{sorry}, no \texttt{native\_decide}):
\texttt{three\_dvd\_gcd\_permanent} proves $3\mid a(n)$ for all $n\ge13$ over
$\mathbb Z/3$, via Lemma A (\texttt{permanent\_eq\_zero\_of\_col\_period}, over
\texttt{Matrix.permanent}) and the $3$-cycle through columns $1,7,13$.
The $p=2$ case can be sharpened past $2\mid a(n)$:

\begin{proposition}\label{prop:four}
$4\mid a(n)$ for all $n\ge4$.
\end{proposition}
\begin{proof}
By the permanent--determinant identity $a(n)=\det(M)+2\Sigma$, where
$\Sigma=\sum_{\sigma\ \mathrm{odd}}\prod_i\gcd(i,\sigma i)$. Smith's theorem gives
$\det(M)=\prod_{k\le n}\varphi(k)$, and $\varphi(3)\varphi(4)=4$ makes
$4\mid\det(M)$ for $n\ge4$. Modulo $2$, $\Sigma\equiv\#\{\sigma\ \mathrm{odd}:
\prod_i\gcd(i,\sigma i)\ \mathrm{odd}\}$. A product $\prod\gcd(i,\sigma i)$ is odd
iff no index has both $i,\sigma(i)$ even, i.e.\ $\sigma$ carries the $\lfloor
n/2\rfloor$ evens into the $\lceil n/2\rceil$ odds; there are exactly
$(\lceil n/2\rceil!)^2$ such $\sigma$ (choose an injection evens$\hookrightarrow$odds,
$\lceil n/2\rceil!/(\lceil n/2\rceil-\lfloor n/2\rfloor)!$ ways, then a bijection
of the rest, $\lceil n/2\rceil!$ ways). The indicator matrix
$G=[\gcd(i,j)\ \mathrm{odd}]=J-ee^\top$ ($e=$ evens) has rank $\le2$, so $\det G=0$
for $n\ge3$; hence the odd and even such $\sigma$ are equinumerous and
$\#\{\sigma\ \mathrm{odd}:\prod\ \mathrm{odd}\}=\tfrac12(\lceil n/2\rceil!)^2$,
even for $n\ge3$. Thus $2\Sigma\equiv0\pmod4$, and $4\mid a(n)$ for $n\ge4$.
\end{proof}
Proposition~\ref{prop:four} is formalized \emph{end to end} in Lean 4/Mathlib
(\texttt{RamaLean/Paper3FourDivides.lean}, standard axioms only, no
\texttt{sorry}): the chain includes Smith's determinant
$\det[\gcd(i,j)]=\prod_{k\le n}\varphi(k)$ (via the factorization $M=LDL^\top$)
and the involution argument for $2\mid\Sigma$ over $\mathbb Z/2$.

\begin{theorem}[the $2$-adic tower]\label{thm:v2}
$v_2(a(n))\to\infty$. Explicitly $v_2(a(n))\ge k+1$ whenever $[1,n]$ contains
$2^k$ integers all of whose prime factors are $\equiv1\pmod{2^k}$; this holds for
all large $n$ by Dirichlet, giving $v_2(a(n))\ge(\tfrac12-o(1))\log_2 n$.
\end{theorem}
\begin{proof}
Fix $k\ge2$. If every prime factor of $m$ is $\equiv1\pmod{2^k}$ then so is every
divisor of $m$, so $\gcd(m,\cdot)\equiv1\pmod{2^k}$: the row of $[\gcd(i,j)]$ at
such an index $m$ is $\equiv(1,\dots,1)\pmod{2^k}$. By Dirichlet there are
infinitely many primes $\equiv1\pmod{2^k}$; for $n$ large pick $2^k$ such
row-indices $R\subseteq[n]$, identify $R$ with $(\mathbb Z/2)^k$, and let
$G\le S_n$ be its regular representation (permuting $R$, fixing $[n]\setminus R$).
Then $|G|=2^k$ and each nontrivial $g\in G$ is a product of $2^{k-1}$ disjoint
transpositions, so $\mathrm{sgn}(g)=(-1)^{2^{k-1}}=1$ (as $k\ge2$), i.e.\ $G\le A_n$.
Writing $a(n)=\det+2\Sigma$ with $\Sigma=\sum_{\mathrm{sgn}\,\sigma=-1}
\prod_i\gcd(\sigma i,i)$, the group $G$ acts on $\{\sigma:\mathrm{sgn}\,\sigma=-1\}$
by $\sigma\mapsto g\sigma$ (sign-preserving since $G\le A_n$; free, being a group
acting by translation). As $g$ permutes $R$ and fixes the rest,
$\prod_i\gcd(g\sigma i,i)\equiv\prod_i\gcd(\sigma i,i)\pmod{2^k}$, so mod $2^k$ the
summand is constant on each $G$-orbit and each orbit (of size $|G|=2^k$) sums to
$0$; hence $2^k\mid\Sigma$. Since $2^{k+1}\mid\det=\prod_{j\le n}\varphi(j)$ for
large $n$ ($v_2(\det)=\sum_{j\le n}v_2\varphi(j)\to\infty$), we get
$2^{k+1}\mid a(n)$, i.e.\ $v_2(a(n))\ge k+1$. The rate follows since the count of
admissible $m\le n$ is $\ge\pi(n;2^k,1)\sim n/(2^{k-1}\log n)$, so $k$ may grow to
$\sim\tfrac12\log_2 n$.
\end{proof}

\begin{remark}
This settles the tower ($k=1$, i.e.\ $2\mid a(n)$, is the involution
$(0\,2)(1\,3)$ of Proposition~\ref{prop:four}; $8\mid a(n)$ for $n\ge17$ uses the
Klein four-group on rows $m=1,5,13,17$, all $\equiv1\bmod4$). The case
$8\mid a(n)$ is \emph{formalized end to end in Lean 4/Mathlib} for $n\ge17$
(theorem \texttt{Paper3Four.eight\_dvd\_permanent}, three standard axioms, no
\texttt{sorry}): the regular Klein-four
$V_4=\langle(c_0c_1)(c_2c_3),\,(c_0c_2)(c_1c_3)\rangle$ acts on the four rows
$m=1,5,13,17$; freeness is immediate (every nontrivial element moves $c_0$), and the
non-disjoint commutation is supplied by the conjugation lemma $e(c\,d)e^{-1}=(ec)(ed)$.
It reuses the orbit-sum engine \texttt{sum\_zmod\_eq\_zero\_of\_free\_invariant}.
Empirically the true
rate is linear, $v_2(a(n))\approx v_2(n!)\sim n$ (with $v_2(n!)-v_2(a(n))\in[-1,6]$ through $n=34$, by
a mod-$2^{128}$ Ryser)---far above the $\tfrac12\log_2 n$ the construction
guarantees, since most terms $\prod_i\gcd(\sigma i,i)$ carry high $2$-adic
valuation while the group action only extracts the ``$\equiv1\bmod 2^k$'' rows.
A \emph{linear} lower rate $v_2(a(n))\gtrsim n/2$ is proved unconditionally in
Theorem~\ref{thm:linear}; this is sharpened to the constant $c=1$,
$v_2(a(n))\ge n-2\lfloor\log_2 n\rfloor-2$, in Theorem~\ref{thm:c1}. The deficit
$v_2(n!)-v_2(a(n))$ appears to be \emph{bounded}: it reaches $6$ at $n=33$ but does not continue
to drift---a factorization $N_0=\sum_v \perm(A_v)\perm(B_v)$ of the dominant $2$-adic grade makes
$v_2(N_0(n))$ computable at $n=65$, giving $v_2(N_0(65))=58$ (deficit $5$), which breaks the
apparent $2,4,6,\dots$ growth along $n=2^k+1$. So conjecturally $v_2(a(n))=v_2(n!)-O(1)$, an even
sharper statement than Theorem~\ref{thm:c1}; a rigorous constant bound remains open.
\end{remark}

\begin{theorem}[linear $2$-adic rate]\label{thm:linear}
For all $n\ge1$,
\[
  v_2\bigl(a(n)\bigr)\ \ge\ \Bigl\lceil \tfrac n2\Bigr\rceil-\lfloor\log_2 n\rfloor-1
  \ =\ \tfrac12 n-O(\log n),
\]
so $v_2(a(n))\to\infty$ \emph{linearly}.
\end{theorem}
\begin{proof}
By Smith's identity $\gcd(i,j)=\sum_{d\mid i,\ d\mid j}\varphi(d)$, expand the permanent one
column at a time:
\[
  a(n)=\sum_{\substack{d=(d_1,\dots,d_n)\\ d_i\mid i\ \forall i}}
        \Bigl(\prod_{i=1}^n\varphi(d_i)\Bigr)\,\perm(M_d),
  \qquad M_d[k,i]=\mathbf 1[\,d_i\mid k\,],
\]
where $\perm(M_d)=\#\{\sigma\in S_n:\ d_i\mid\sigma(i)\ \forall i\}$. Fix a tuple $d$ and put
$c_1=\#\{i:d_i=1\}$, $c_2=\#\{i:d_i=2\}$, $c_3=\#\{i:d_i\ge3\}$. The $c_1$ columns with $d_i=1$
are all-ones, so in any admissible $\sigma$ they take, in any order, the $c_1$ values left free by
the other columns; hence $\perm(M_d)=c_1!\cdot N$ for some integer $N$, and
$v_2(\perm(M_d))\ge v_2(c_1!)=c_1-s_2(c_1)$ (where $s_2$ is the binary digit sum). Since
$\varphi(d)$ is even for every $d\ge3$, we have $v_2\!\bigl(\prod_i\varphi(d_i)\bigr)\ge c_3$.
Finally $d_i=2$ forces $2\mid i$, so $c_2\le\lfloor n/2\rfloor$ and
$c_1+c_3=n-c_2\ge\lceil n/2\rceil$. Thus every term satisfies
\[
  v_2\Bigl(\textstyle\prod_i\varphi(d_i)\,\perm(M_d)\Bigr)
  \ \ge\ (c_1-s_2(c_1))+c_3
  \ =\ (c_1+c_3)-s_2(c_1)
  \ \ge\ \Bigl\lceil\tfrac n2\Bigr\rceil-s_2(c_1)
  \ \ge\ \Bigl\lceil\tfrac n2\Bigr\rceil-\lfloor\log_2 n\rfloor-1,
\]
using $s_2(c_1)\le\lfloor\log_2 c_1\rfloor+1\le\lfloor\log_2 n\rfloor+1$. As $a(n)$ is an integer
combination of these terms, $v_2(a(n))$ is at least this common bound.
\end{proof}

\begin{remark}
The method is tight in \emph{shape} (linear) but not in constant: numerically $v_2(a(n))\approx
v_2(n!)\sim n$ (slope $1$; deficit $v_2(n!)-v_2(a(n))\in[-1,6]$ through $n=34$, and $D(2^k+1)=2k-4$ at the peaks $n=2^k+1$ for
$k\le5$, with the corresponding statement for the weight-zero grade proved at every $k$ of odd
achiever parity; see Section~\ref{sec:deficit}), and the stronger
$v_2(a(n))\ge v_2(\lceil n/2\rceil!)$ holds for every $n\le20$. The factor-$2$ loss comes from
using only the $\lceil n/2\rceil$ odd positions' freedom; the even positions carry a comparable
$2$-adic contribution that a single application of the all-ones lemma does not see.
\end{remark}

\begin{remark}[formalization]
Theorem~\ref{thm:linear} is \emph{formalized end to end in Lean 4/Mathlib}
(\texttt{RamaLean/Paper3LinearRate.lean}, theorem \texttt{two\_pow\_dvd\_permanent}:
$2^{\lceil n/2\rceil-\lfloor\log_2 n\rfloor-1}\mid a(n)$, three standard axioms, no
\texttt{sorry})---the \emph{same} constant $c=\tfrac12$ proved above. The Lean proof formalizes
the permanent's Smith expansion $a(n)=\sum_x\bigl(\prod_i\varphi(x_i)\bigr)\perm(M_x)$
(column multilinearity via \texttt{Finset.prod\_univ\_sum}), the all-ones-column divisibility
$c_1!\mid\perm(M_x)$ (a free \texttt{Perm(Fin }$c_1$\texttt{)} action, from the reusable engine
\texttt{card\_dvd\_sum\_of\_free\_invariant}: a free finite-group action with invariant summand
makes $|G|$ divide the sum over any commutative ring), the \emph{exact} factorial valuation
$v_2(c_1!)=c_1-s_2(c_1)$ (\texttt{Nat.sub\_one\_mul\_factorization\_factorial}), and the
$2$-adic bookkeeping.
\end{remark}

\begin{theorem}[the growth constant: lower bound $c=1$]\label{thm:c1}
For all $n\ge1$,\quad $v_2\bigl(a(n)\bigr)\ \ge\ n-2\lfloor\log_2 n\rfloor-2$,\quad hence
$\liminf_n v_2(a(n))/n\ge1$.
\end{theorem}
\begin{proof}
Write $\gcd(i,j)=2^{\min(v_2 i,\,v_2 j)}h(i,j)$ with $h(i,j)=\gcd(\mathrm{odd}(i),\mathrm{odd}(j))$
\emph{odd}, so $a(n)=\sum_\sigma 2^{w(\sigma)}\prod_i h(\sigma i,i)$ with
$w(\sigma)=\sum_i\min(v_2(\sigma i),v_2(i))$. Only pairs with both indices even contribute to $w$
(an odd index has $v_2=0$). Grouping the permutations by their even$\to$even sub-injection
$\rho$ (an even position sent to an even value),
\[
  a(n)=\sum_\rho 2^{w_\rho}\Bigl(\prod_{i\in\mathrm{dom}\,\rho}h(\rho i,i)\Bigr)\mathrm{Rest}_\rho,
  \qquad w_\rho=\!\!\sum_{i\in\mathrm{dom}\,\rho}\!\!\min(v_2(\rho i),v_2(i))\ \ge\ |\rho|,
\]
where $\mathrm{Rest}_\rho=\perm(H_\rho)$ and $H_\rho$ is the $(n-|\rho|)\times(n-|\rho|)$ matrix
equal to $h$ off the (remaining even)$\times$(remaining even) corner and $0$ on that corner
(remaining even positions cannot reuse even values). As $\prod h$ is odd,
$v_2(\text{term}_\rho)=w_\rho+v_2(\mathrm{Rest}_\rho)$.

\emph{Zeroed-corner lemma:} if $M$ is $s\times s$, odd off a $b\times b$ corner and $0$ on it,
then $v_2(\perm M)\ge s-2\lfloor\log_2 s\rfloor-2$. Indeed the $b$ corner rows are forced into
the $s-b$ non-corner columns, so
$\perm M=\sum_{|T|=b}\perm(B[\cdot,T])\,\perm([D\mid E[\cdot,C_2\!\setminus\!T]])$, a sum of
products of two \emph{all-odd} permanents of sizes $b$ and $s-b$; each all-odd $k\times k$
permanent has $v_2\ge k-\lfloor\log_2 k\rfloor-1$ (Theorem~\ref{thm:linear}'s engine: a
permanent of an odd matrix $=J+2E$ expands multilinearly as $\sum_T 2^{|T|}\perm_T$ with
$(k-|T|)!\mid\perm_T$), so each summand has $v_2\ge s-2\lfloor\log_2 s\rfloor-2$.

Therefore $v_2(\mathrm{Rest}_\rho)\ge(n-|\rho|)-2\lfloor\log_2 n\rfloor-2$ and
$v_2(\text{term}_\rho)\ge|\rho|+(n-|\rho|)-2\lfloor\log_2 n\rfloor-2=n-2\lfloor\log_2 n\rfloor-2$;
as $a(n)$ is their sum, $v_2(a(n))\ge n-2\lfloor\log_2 n\rfloor-2$.
\end{proof}
\begin{remark}[formalization --- fully machine-checked]
Theorem~\ref{thm:c1} is \emph{formalized end to end in Lean\,4/Mathlib}:
\texttt{Paper3C1.two\_pow\_dvd\_permanent\_c1}:
$2^{\,n-2\lfloor\log_2 n\rfloor-2}\mid(\mathtt{gcdMat}\,n).\mathtt{permanent}$, standard axioms
\texttt{[propext, Classical.choice, Quot.sound]}, no \texttt{sorry}. The Lean proof is even
\emph{shorter} than the hand proof above: rather than reindex $a(n)$ by permutation grade, it
observes that the gcd matrix \emph{is} a zeroed-corner matrix---$\gcd(k{+}1,i{+}1)$ is even iff
both $k{+}1,i{+}1$ are even, so $a(n)=\perm\!\bigl(2\,ee+[\neg(k,i\text{ both even})]\bigr)$---and
applies the reusable engine \texttt{ZeroedCorner.two\_pow\_dvd\_permanent\_zeroed}. That engine
(column-multilinear expansion, each term carrying two groups of identical mask-columns) rests on a
\emph{two-group factorial} \texttt{TwoGroup.two\_factorial\_dvd\_permanent}
($|A|!\,|B|!\mid\perm M$ for two disjoint groups of identical rows, via a free
$\mathrm{Sym}(A)\times\mathrm{Sym}(B)$ action---the product-group generalisation of the
all-ones-rows kernel behind Theorem~\ref{thm:linear}). The constant $c=1$ is optimal
\emph{conditional} on the empirical upper bound $v_2(a(n))\le v_2(n!)+O(1)\sim n$ (observed through
$n=34$, unproven); what is established here is the lower bound, hence $\liminf_n v_2(a(n))/n\ge1$.
\end{remark}

\section{Growth rate}
Write $b(n)=a(n)/n!=\mathbb E_\sigma\prod_i\gcd(i,\sigma(i))$.

\begin{theorem}\label{thm:up}
$(a(n)/n!)^{1/n}=O\big((\log n)^{\log 2}\big)$; equivalently
$a(n)\le n!\,(\log n)^{(\log 2+o(1))n}$.
\end{theorem}
\begin{proof}
$\perm(M)\le\prod_i R_i$ with $R_i=\sum_j\gcd(i,j)\le n f(i)$,
$f(i)=\sum_{d\mid i}\varphi(d)/d=\prod_{p^a\|i}(1+a(1-1/p))$. By Stirling,
$(a/n!)^{1/n}\le e\exp\big(\tfrac1n\sum_{i\le n}\log f(i)\big)$. Swapping sums,
$\tfrac1n\sum_{i\le n}\log f(i)=\sum_{p\le n}g(p)+o(1)$ with
$g(p)=\sum_{k\ge1}(1-\tfrac1p)p^{-k}\log(1+k(1-\tfrac1p))=\tfrac{\log2}{p}
+O(p^{-2})$. By Mertens, $\sum_{p\le n}g(p)=(\log2)\log\log n+O(1)$, giving the
bound.
\end{proof}

\begin{lemma}[entropy lower bound]\label{lem:ent}
For every doubly stochastic $w=(w_{ij})$,
\[
\frac{a(n)}{n!}\ge \frac{1}{n^n}\exp\!\Big(\sum_{ij}w_{ij}\log\gcd(i,j)+H(w)\Big),
\qquad H(w)=-\sum_{ij}w_{ij}\log w_{ij}.
\]
\end{lemma}
\begin{proof}
Let $\mathrm{cap}(M)=\inf_{\prod z_j=1}\prod_i(Mz)_i$. Weighted AM--GM (using
$\sum_j w_{ij}=1$) gives $(Mz)_i=\sum_j M_{ij}z_j\ge\prod_j (M_{ij}z_j/w_{ij})
^{w_{ij}}$, so, using $\sum_i w_{ij}=1$ and $\prod_j z_j=1$,
$\prod_i(Mz)_i\ge\prod_{ij}(M_{ij}/w_{ij})^{w_{ij}}=\exp(\sum w_{ij}\log M_{ij}
+H(w))$. Hence $\mathrm{cap}(M)\ge\exp(\sum w\log M+H(w))$, and Gurvits'
theorem $\perm(M)\ge(n!/n^n)\,\mathrm{cap}(M)$ (which for doubly stochastic
matrices is exactly Van der Waerden--Falikman--Egorychev) finishes.
\end{proof}

\begin{theorem}[growth rate]\label{thm:c}
$(a(n)/n!)^{1/n}\to\infty$. Writing $(a(n)/n!)^{1/n}=(\log n)^{\theta_n}$, the
exponent obeys $\liminf_n\theta_n\ge2\ln\varphi-\varphi^{-2}=0.5804\ldots$ and
$\limsup_n\theta_n\le\log2$; numerically $\theta_n\approx0.68$ but in a strongly
pre-asymptotic regime ($\log\log n\approx2$ even at $n=3000$), so the exact
value---in particular whether $\theta=\log2$---is open.
\end{theorem}
\begin{proof}
Theorem~\ref{thm:up} gives $\theta_n\le\log2+o(1)$; we prove the lower bound
$\theta_n\ge2\ln\varphi-\varphi^{-2}-o(1)$ (in particular $(a/n!)^{1/n}\to
\infty$). Fix a cutoff $P$. Since
$\prod_{p\le P,\,p\mid i,\,p\mid j}p$ divides $\gcd(i,j)$, monotonicity of the
permanent gives $a(n)\ge\perm(\tilde M)$ with
$\log\tilde M_{ij}=\sum_{p\le P}(\log p)\,\mathbf1[p\mid i]\,\mathbf1[p\mid j]$.
Partition $[1,n]$ by divisibility pattern $S(i)=\{p\le P:p\mid i\}$: for
$S\subseteq\{p\le P\}$, whenever $\prod_{p\le P}p\mid n$ the interval is a union
of complete residue systems mod $\prod_{p\le P}p$, so
\[
 |V_S|=n\mu_S\ \text{\emph{exactly}},\qquad
 \mu_S=\prod_{p\in S}\tfrac1p\prod_{p\notin S}\big(1-\tfrac1p\big)
\]
(for general $n$, $|V_S|=n\mu_S+O(2^{\pi(P)})$, the error constant in $n$). Let
$\rho=\bigotimes_{p\le P}\rho^{(p)}$ be the product of the optimal Bernoulli
couplings of Theorem~\ref{thm:gr}, and set
$w_{ij}=n\,\rho_{S(i)S(j)}/(|V_{S(i)}||V_{S(j)}|)$, constant on the
$2^{\pi(P)}\times2^{\pi(P)}$ blocks. When $\prod_{p\le P}p\mid n$ this $w$ is
\emph{exactly} doubly stochastic, and a direct block computation gives
\[
 \tfrac1nF(w)-\log n=\sum_{p\le P}\Big[(\log p)\,\rho^{(p)}_{11}
 -D\big(\rho^{(p)}\,\|\,\mathrm{Bern}(\tfrac1p)^{\otimes2}\big)\Big]
 =\sum_{p\le P}g_p .
\]
By Lemma~\ref{lem:ent}, $(a(n)/n!)^{1/n}\ge e^{\sum_{p\le P}g_p}$ along
$\prod_{p\le P}p\mid n$ (for general $n$ a Sinkhorn correction on the
$2^{\pi(P)}$ blocks changes $F$ by $O_P(1)=o(n)$, giving the same limit); hence
$\liminf_n(a(n)/n!)^{1/n}\ge e^{\sum_{p\le P}g_p}$ for \emph{every} $P$. By
Theorem~\ref{thm:gr} $\sum_{p\le P}g_p\to\infty$ as $P\to\infty$, so
$(a(n)/n!)^{1/n}\to\infty$.
\end{proof}

\begin{remark}[{the exponent gap $[\,2\ln\varphi-\varphi^{-2},\ \log2\,]$}]
Since $\perm(M)\le\mathrm{cap}(M)\le\prod_iR_i$ and $\perm(M)\ge(n!/n^n)\,
\mathrm{cap}(M)$, one has $(a(n)/n!)^{1/n}=\Theta\big(e^{\,\max_wF(w)/n-\log n}\big)$
(constant in $[1,e]$), so the exponent is governed by the true capacity
$\log\mathrm{cap}(M)=\max_wF(w)$. Our construction couples primes
\emph{independently}, reaching only $2\ln\varphi-\varphi^{-2}$; the bound
$\theta\le\log2$ uses $R_i=\sum_j\gcd(i,j)$, which feels all primes at once.
Numerically the capacity exponent rises---$\theta_n=0.672,0.685,0.689$ at
$n=10^2,10^3,3\!\cdot\!10^3$, with $\log2-\theta_n=0.021,0.008,0.004$ apparently
$\to0$---but this is unreliable: the \emph{upper} rate $E_i[\log f(i)]/\log\log n$
itself converges to $\log2$ only \emph{logarithmically slowly} (still $\approx0.74$
at $n=10^7$), so at any feasible size the data cannot separate $\theta=\log2$
from a strictly smaller constant (the Sinkhorn deficit
$\sum_i\log R_i-\log\mathrm{cap}(M)$ is $\approx0.026\,n\log\log n$ at these $n$,
consistent with either). The
obstruction to a clean answer is the hard constraint $\sum_pv_p(i)\log p=\log i$,
which couples all primes and defeats the independent-prime (product) analysis:
the golden-ratio value $2\ln\varphi-\varphi^{-2}$ is the exponent of that
\emph{mean-field} strategy, cross-prime correlation provably beats it, and how
far it goes---to $\log2$ or short of it---is the open question.
\end{remark}

\begin{theorem}[per-prime gain; a golden-ratio constant]\label{thm:gr}
For a prime $p$ define the divisibility gain
\[
 g_p=\sup_{\rho}\Big[(\log p)\,\rho_{11}
 -D\big(\rho\,\|\,\mathrm{Bern}(\tfrac1p)^{\otimes2}\big)\Big],
\]
the supremum over $2\times2$ couplings $\rho$ with both marginals
$\mathrm{Bern}(1/p)$---the law of $\mathbf1[p\mid i]$---of the divisibility
indicator with itself. Then $g_p>0$ for every $p$, and
\[
 \lim_{p\to\infty}p\,g_p=2\ln\varphi-\varphi^{-2}=2\ln\varphi+\varphi-2
 =0.5804576\ldots,\qquad\varphi=\tfrac{1+\sqrt5}2,
\]
attained at $\rho_{11}=t^{*}/p$ with $t^{*}=\varphi^{-2}=\tfrac{3-\sqrt5}2$, the
golden section (root of $(1-t)^2=t$). Hence $\sum_pg_p=\infty$, with partial sums
$\sim(2\ln\varphi-\varphi^{-2})\log\log P$.
\end{theorem}
\begin{proof}
Write $\varepsilon=1/p$ and $\rho_{11}=\varepsilon y$ ($y\in[0,1]$), so
$\rho_{10}=\rho_{01}=\varepsilon(1-y)$ and $\rho_{00}=1-2\varepsilon+\varepsilon
y$. In the four relative-entropy terms the only unbounded piece,
$\rho_{11}\log(\rho_{11}/\varepsilon^2)=\varepsilon y\log y+\varepsilon y\log p$,
has its $\varepsilon y\log p$ cancelled \emph{exactly} by the gain
$(\log p)\rho_{11}$; the rest expands with an absolute $O(\varepsilon^2)$ error
(no logarithm), giving
\[
 g_p(y)=\varepsilon\,\Psi(y)+O(\varepsilon^2),\qquad
 \Psi(y)=-y-y\log y-2(1-y)\log(1-y).
\]
As $p\varepsilon=1$, $p\,g_p(y)=\Psi(y)+O(1/p)$; maximizing $\Psi$ at
$\Psi'(y)=\log\frac{(1-y)^2}{y}=0$, i.e.\ $(1-y)^2=y$ and $y=\varphi^{-2}$,
yields $p\,g_p\ge\Psi(\varphi^{-2})-O(1/p)=2\ln\varphi-\varphi^{-2}-O(1/p)$,
whence $\liminf_p p\,g_p\ge2\ln\varphi-\varphi^{-2}$ (numerically an equality,
$p\,g_p\uparrow0.58046$). Positivity at finite $p$ is the strict suboptimality
of the independent coupling. The explicit $O(1/p)$ makes
$g_p\ge(2\ln\varphi-\varphi^{-2}-O(1/p))/p$, so $\sum_pg_p=\infty$.
\end{proof}
\begin{remark}[the finer coupling]
Using the full valuation $v_p(i)$ in place of $\mathbf1[p\mid i]$ gives
$G_p=\sup_\pi[(\log p)\E_\pi\min(a,b)-D(\pi\|q_p\otimes q_p)]\ge g_p$
($q_{p,k}=(1-\tfrac1p)p^{-k}$), with the clean identity
$\E_{q_p\otimes q_p}[p^{\min}]=1+\tfrac1p$---so $G_p=\log(1+\tfrac1p)
-\mathrm{SB}_p$, a Schr\"odinger-bridge cost---and the \emph{same} limit
$pG_p\to2\ln\varphi-\varphi^{-2}$. The coarse divisibility gain $g_p$ already
suffices for Theorem~\ref{thm:c}. The divergence is genuinely \emph{annealed}
and multi-scale---invisible to the uniform measure (Jensen sees only $0.570$),
to any $O(1)$-tilt, and to any single block---and the golden section
$\varphi^{-2}$ is the signature of the entropy tradeoff $\Psi$.
\end{remark}

\section{The 2-adic deficit and two odd-permanent engines}\label{sec:deficit}

Theorem~\ref{thm:c1} gives $v_2(a(n))\ge n-2\lfloor\log_2n\rfloor-2$, i.e.\ the deficit
$D(n):=v_2(n!)-v_2(a(n))$ is $O(\log n)$. We now explain its \emph{value}. Exact Gray-code Ryser data
at the peaks $n=2^k+1$ gives $v_2(a)=3,5,11,25$ for $k=2,3,4,5$, which is
\begin{equation}\label{eq:defline}
   v_2\big(a(2^k+1)\big)=2^k-2k+3,\qquad D(2^k+1)=2k-4\sim2\log_2 n
   \qquad(k=2,3,4,5).
\end{equation}
We stress that \eqref{eq:defline} is data, not a theorem: what follows explains those four values
and proves the corresponding statement for the weight-zero grade $N_0$ of $a$, at every $k$ with odd
achiever parity. Whether $v_2(a)=v_2(N_0)$ persists beyond $k=5$ we do not know.
The mechanism behind~\eqref{eq:defline} is two exact 2-adic valuation theorems for odd permanents.
For an $s\times s$ all-odd matrix write $M=2e+J$ ($J$ all-ones); the column-multilinear expansion is
$\perm(2e+J)=\sum_{k=0}^{s}2^k(s-k)!\,\pi_k(e)$, $\pi_k=\sum_{|R|=|S|=k}\perm(e[R,S])$, so term $k$ has
$v_2=s-s_2(s-k)+v_2(\pi_k)$ (with $s_2$ the binary digit sum).

\begin{theorem}[$A$-engine]\label{thm:engineA}
If $s=2^a$, $a\ge2$, and $\pi_1(e)=\sum_{i,j}e_{ij}$ is odd, then $v_2(\perm(2e+J))=2^a-a$.
\end{theorem}
\begin{proof}
The $k=1$ term has $v_2=2^a-s_2(2^a-1)+v_2(\pi_1)=2^a-a$ (as $\pi_1$ is odd); the $k=0$ term has
$v_2=2^a-1>2^a-a$; and for $2\le k\le2^a$ every $2^a-k\in[0,2^a-2]$ has $s_2\le a-1$ (only $2^a-1$ has
$a$ ones), so those terms have $v_2\ge2^a-a+1$. The $k=1$ term is the unique minimum: no cancellation,
hence the claim.
\end{proof}

\begin{theorem}[$B$-engine]\label{thm:engineB}
If $s=2^a+1$, $a\ge3$, and $\pi_2(e)$ is odd, then $v_2(\perm(2e+J))=2^a+1-a$.
\end{theorem}
\begin{proof}
Now $\max_k s_2(s-k)=a$ is attained uniquely at $k=2$ ($s-k=2^a-1$), giving that term
$v_2=2^a+1-a+v_2(\pi_2)=2^a+1-a$ (as $\pi_2$ is odd); all other $k$ have $s_2(s-k)\le a-1$, hence
$v_2\ge2^a+2-a$. Unique minimum, no cancellation.
\end{proof}
Theorem~\ref{thm:engineA} is the tight matching upper bound to Theorem~\ref{thm:linear}, and it is
\textbf{fully machine-checked}: the lower half $2^{2^a-a}\mid\perm(2e+J)$
(\texttt{OddEngine.two\_pow\_dvd\_permanent\_odd\_pow2}) and the upper half
$2^{2^a-a+1}\nmid\perm(2e+J)$ under $\pi_1$ odd (\texttt{OddEngine.engine\_upper}), together giving
$v_2(\perm(2e+J))=2^a-a$ exactly (Lean~4/Mathlib, standard axioms, no \texttt{sorry}). The formalized
proof splits the column-subset expansion by $|t|$: the $|t|{=}0$ term is $\perm J=s!$, the $|t|{\ge}2$
terms vanish mod $2^{2^a-a+1}$ (via the binary-complement identity $s_2(x)+s_2(2^a{-}1{-}x)=a$), and the
$|t|{=}1$ terms sum to $2(s{-}1)!\,\pi_1$ (via the orbit--stabilizer fiber count
$\#\{\sigma:\sigma j=r\}=(s{-}1)!$).

For odd $n$ the weight-$0$ grade of $a(n)=\sum_w2^wN_w$ factors as $N_0=\sum_v\perm(A_v)\perm(B_v)$,
where $A_v$ (size $\lfloor n/2\rfloor$) and $B_v$ (size $\lfloor n/2\rfloor+1$) are all-odd. At the
peaks $\lfloor n/2\rfloor=2^{k-1}$, so Theorems~\ref{thm:engineA}--\ref{thm:engineB} apply and, when
$\pi_1(A_v),\pi_2(B_v)$ are both odd,
\[
   v_2(\perm A_v)+v_2(\perm B_v)=\big(2^{k-1}-(k-1)\big)+\big(2^{k-1}+1-(k-1)\big)=2^k-2k+3 .
\]
These are the minimum-valuation terms of $N_0$, so the leading-term rule gives
\[
   v_2(N_0)=2^k-2k+3\iff\#\{v:\pi_1(A_v),\pi_2(B_v)\text{ both odd}\}\text{ is odd},
\]
reproducing $v_2(N_0)=11,25,58$ at $k=4,5,6$ exactly (the $58$ is the even-parity lift off $55$). This
achiever count is computable in $O(n\log n)$ with no permanent (a divisor sieve); it is odd for
$k\in\{4,5,7,8,11,14,16,17,18,22,23\}$ within $k\le23$ ($11$ of $20$, recurring), so
$D_{N_0}(2^k+1)=2k-4$ there, reaching $42$ at $k=23$ and setting a new record at every odd-parity peak.

\begin{conjecture}\label{conj:io}
The achiever count is odd for infinitely many $k$; equivalently $D_{N_0}(2^k+1)$ is unbounded,
$=\Theta(\log n)$ along the peaks.
\end{conjecture}
For the full permanent one needs additionally $v_2(a)=v_2(N_0)$, which holds at $k=2,3,4,5$ (though the
grade-margins $0,1,3,1$ are small---a bare tie at $k=2$---so this dominance is fragile).

\paragraph{An empirical refinement.} Computing $C(k)$ to $k=35$ and stratifying by $k\bmod 6$ reveals
one residue class that is not pseudorandom: $C(k)=1$ for \emph{every} $k\equiv5\ (\mathrm{mod}\ 6)$
tested ($k=5,11,17,23,29,35$, i.e.\ $6/6$), while each of the other five classes is genuinely mixed.
This suggests the sharper
\begin{conjecture}\label{conj:mod6}
$C(k)=1$ for all $k\equiv5\pmod6$; equivalently $D_{N_0}(2^k+1)=2k-4$ on this class. This would imply
Conjecture~\ref{conj:io}.
\end{conjecture}
\noindent This is a lead, not a near-theorem: six independent structural approaches (dominant-stratum, achiever involution, small-prime CRT localization,
period-$6$ recurrence, $\chi_4$-reduction, and a MacWilliams \cite{MS} character-sum expansion) each terminate at the same parity barrier.
The achiever counts on the class are $1,211,12659,\dots$ (odd but large and growing), so the
constancy is the parity of a genuine correlation, not a lone surviving term. Two facts survive as
rigorous byproducts. First, a cleaner form of the target: the achiever set is exactly the
intersection of \emph{two affine $\mathbb F_2$ hyperplanes} in the $\equiv3\ (4)$ prime-power divisor
profile of $v$, $A=\{v\ \mathrm{odd}\le n:R(v)=1\wedge G(v)=1\}$ with both $R,G$ affine-linear (the
apparent quadratic $B$-side terms cancel since $[D\mid v][E\mid v]=[\operatorname{lcm}(D,E)\mid v]$).
Second, $C$ is a genuine \emph{cross-stratum} quantity, not a localized invariant: at $k=23$ the
bilinear correlation term and its diagonal and off-diagonal parts all vanish while $C=1$ is carried
by the linear term---so no single stratum drives the class, and any ``dominant term'' argument fails
at $k=23$. The class is distinguished from the mixed class $k\equiv1\ (6)$ only by $2^k\bmod7$ ($4$
versus $2$; both give $3\parallel n$), but that digit acts through the full $B$-side correlation, not
as a lever. Proving Conjecture~\ref{conj:mod6} therefore appears to require the same
correlation-parity control as Conjecture~\ref{conj:io}: it sharpens the target to a clean residue
class and a cleaner (two-hyperplane) form, without lowering the barrier. The
remaining analytic handle---the $\mathbb F_2$-Fourier expansion $|A|=\tfrac14(N-S_R-S_G+S_{R\oplus
G})$, whose cross-term $S_{R\oplus G}=\sum_{v\ \mathrm{odd}\le n}f(v)$ has $f$ multiplicative---meets
the same barrier: $f$ is a genuinely $k$-dependent multiplicative $\pm1$ function (not a
Dirichlet character: its $f(p)=-1$ density swings $23\%\!\to\!79\%\!\to\!19\%$ across consecutive
$k$), so no closed form or class-number identity applies; and since $N\equiv1\ (16)$, the criterion
$|A|$ odd $\iff N-S_R-S_G+S_{R\oplus G}\equiv4\ (8)$ is \emph{identical} to the conjecture, not a
reduction. Two structural correlates of the class survive computation ($k\le17$): $3\mid S_{R\oplus G}$ (equivalently $S_G\equiv1,\ S_{R\oplus G}\equiv0\bmod3$, both constant on the class) on every class member and no non-member ($9/9$), and the flip-set orientation is uniformly sparse
($c_k=0$) on $k\equiv5\ (6)$ whereas $k\equiv1\ (6)$ flips regime between members. Neither is a
parity lever ($\mathrm{mod}\,3\perp\mathrm{mod}\,8$), so the honest status is: \emph{empirically
true, analytically barriered}. The root obstruction is precise: $|A|$ is not the size of an affine
$\mathbb F_2$-subspace (which would be a power of $2$), because $v$ ranges over \emph{integers} $\le
n$, not over $\mathbb F_2^{|P|}$---so $|A|$ is a \emph{lattice-point count} in an affine variety, with
$|A|/(N/4)$ wandering in $[0.24,0.96]$ (measured $k\le17$), the anti-correlation being exactly the
growing character sums. The domain constraint breaks the $\mathbb F_2$-linear structure, which is why it cannot force the parity. None of the standard analytic techniques (exponential sums / circle method, automatic-sequence and sum-of-digits methods, Dirichlet-series / transfer-operator analysis of the doubling map, and the random-multiplicative model) applies: each controls magnitude or equidistribution, whereas the target is a $2$-adic last bit. The parity of $|A|$ thus sits in the \emph{Selberg parity-problem} family (two Liouville-type multiplicative $\pm1$ functions), entangled with a Gelfond-type digit correlation and modelled by random-multiplicative fair-coin parity---the canonical warning being $p(n)\bmod2$, where an exact circle-method formula yields nothing and only modular/$\ell$-adic methods (absent here) succeed. The conjecture is genuinely open, and expected to require a new exact-arithmetic, not analytic, input.

\paragraph{The arithmetic of the parity: $\chi_4$ and the orders of $2$.} The parity in
Conjecture~\ref{conj:io} is not opaque---it has an exact arithmetic form that identifies precisely why
it is hard. Since $\gcd$ of two odd numbers is odd and $(g-1)/2$ is odd iff $g\equiv3\ (4)$, the parity
$\pi_1(A_v)\bmod2$ counts pairs whose $\gcd\equiv3\ (4)$:
\[
   \pi_1(A_v)\equiv G_3+g_3(v)\pmod2,\qquad
   g_3(v):=\#\{\,i\le 2^{k-1}:\gcd(i,v)\equiv3\ (4)\,\},
\]
with $G_3$ the same count over all values. Writing $[\,n\equiv3\ (4)\,]=(1-\chi_4(n))/2$ for the
non-principal character $\chi_4$ mod $4$ and Dirichlet-inverting, only prime-power divisors survive
modulo $2$:
\begin{equation}\label{eq:chi4}
   g_3(v)\ \equiv\ \sum_{\substack{q\equiv3\ (4)\ \mathrm{prime}\\ a\ge1,\ q^a\mid v}}
      \Big\lfloor \tfrac{2^{k-1}}{q^a}\Big\rfloor \pmod 2 ,
\end{equation}
and $\lfloor 2^{k-1}/q^a\rfloor\bmod2$ is exactly the $(k{-}1)$-th \emph{binary digit of $1/q^a$},
periodic in $k$ with period $\mathrm{ord}_{q^a}(2)$. Identity~\eqref{eq:chi4} is in fact a
\emph{theorem}: for odd $g$, reducing the factorization mod $4$ gives
$g\equiv(-1)^{\sum_{q\equiv3(4)}v_q(g)}\ (4)$, so $[\,g\equiv3\ (4)\,]$ is the parity of
$\sum_{q\equiv3(4)}v_q(g)=\sum_{q^a}[q^a\mid g]$, and summing over entries yields~\eqref{eq:chi4};
$\pi_2(B_v)$ obeys the same skeleton with different digit weights. Thus the
achiever count $C(k)=\sum_v P_1(v)P_2(v)$ is a \emph{correlation, over all $v$, of two parities built
from the multiplicative orders $\mathrm{ord}_q(2)$ of the primes $q\equiv3\ (4)$ up to $2^k$}. As
$k\to\infty$ new such primes continually enter, each contributing a bit governed by its own order; the
empirical pseudorandomness of $C(k)\bmod2$ is the (conjectural) pseudorandomness of these orders. This
places Conjecture~\ref{conj:io} against \emph{two} known walls: the distribution of $\mathrm{ord}_q(2)$
is Artin-primitive-root territory (unconditional only in part; sharp forms need GRH), and extracting a
\emph{parity} of a counting function from density information is the classical parity barrier. So the
deficit's unboundedness is not merely ``open''---it reduces to, and is obstructed by, open problems on
the orders of $2$.

\paragraph{A handle beyond the character-sum barrier: transcendence.} The parity barrier is a barrier
only in the analytic frame, where one controls the \emph{size} $|\sum\chi|$ of a character sum but
never the \emph{parity} of the underlying count. Passing to the automata/transcendence frame removes
exactly that blindness. The parity sequence $C(k)$ is a sequence over $\mathbb F_2$, and by the proven
equivalence $D_{N_0}(2^k+1)=2k-4\iff C(k)=1$ we have a clean dichotomy.
\begin{theorem}[reformulation]\label{thm:transc}
Let $F(x)=\sum_k C(k)\,x^k\in\mathbb F_2[[x]]$. Then $D_{N_0}(2^k+1)$ is unbounded iff $C(k)$ is not
eventually $0$ iff $F$ is not a polynomial. Moreover each of the following implies unboundedness, in
increasing strength: $F$ is transcendental over $\mathbb F_2(x)$ $\Rightarrow$ $C$ is not
$2$-automatic $\Rightarrow$ $C$ is not eventually periodic $\Rightarrow$ unbounded.
\end{theorem}
\noindent The middle link is Christol's theorem ($C$ is $2$-automatic iff $F$ is algebraic over
$\mathbb F_2(x)$); the outer links are that eventually-zero $\Rightarrow$ eventually-periodic
$\Rightarrow$ automatic. Unlike a size bound, algebraicity \emph{is} sensitive to parity, so this is a
genuine escape route rather than a restatement. It also has an explicit pole obstruction. Since each
word is periodic with period $m_v=\mathrm{ord}_v(2)$ once $v$ enters at scale $\kappa_v$,
\[
   F(x)\ =\ \sum_v x^{\kappa_v}\,\frac{A_v(x)}{1-x^{m_v}}\pmod 2,\qquad \deg A_v<m_v ,
\]
a legitimate power series (each coefficient is a finite sum). A rational function has finitely many
poles, but the denominators place poles at $m_v$-th roots of unity with $m_v$ unbounded; so $F$
rational forces the combined residue at every primitive $d$-th root of unity to vanish for all large
$d$. By Zsygmondy each order $d\neq6$ carries a primitive prime $q_d\mid2^d-1$ injecting a pole at
$\zeta_d$ with residue $\propto A_{q_d}(\zeta_d)$---and $A_{q_d}$ is, by the Prime-Word Theorem below,
the binary-expansion word of $1/q_d$. Rationality thus demands an explicit infinite family of
cancellations among the reciprocal-digit words of Mersenne primitive primes---a concrete no-conspiracy
condition, though not obviously unattainable, since the smallest prime of a given
order need not be a Mersenne prime.
A divisor-sieve from the linearity lemma computes $C(k)$ in $O(n\log\log n)$ (validated against the
direct engine, $k\le22$), giving $C(k)$ for $k\le36$: density $\tfrac{18}{33}$ and a
Berlekamp--Massey linear-complexity profile $1,2,5,6,10,13,16$ that climbs in plateau-then-jump
fashion (it held at $13$ for six terms, jumped to $16$ at $k=32$, and has held there since). These are
\emph{computed}, not proved: the full-$C(k)$ sieve is archived (\texttt{code/rust\_grade/fastC.rs}),
but the Berlekamp--Massey step is not archived in this release, so we draw no rigorous
non-periodicity conclusion from the complexity profile---and $33$ terms cannot by themselves
separate a high-complexity rational $F$ from a transcendental one. Crucially, either way unboundedness
needs only that $F$ is not a \emph{polynomial}, i.e.\ $C(k)$ is not eventually zero---and the direct
evidence for that is strong and independent of the transcendence question: $C(k)=1$ at $15$ of the
$29$ peaks $k\le32$, so $D_{N_0}(2^k+1)=2k-4$ there; the six largest such deficits are
$40,42,44,54,56,58$ at $k=22,23,24,29,30,31$, with $C(k)=0$ in the gap $25\le k\le28$. What Theorem~\ref{thm:transc} contributes is not a resolution but a relocation:
the parity in question is a coefficient of an explicit $\mathbb F_2$ power series, an object to which
the $2$-kernel and Mahler's method apply and to which the analytic parity barrier does not speak.

\paragraph{Adaptive peaks: weakening the conjecture.} The parity barrier above blocks the \emph{fixed}
sequence $n=2^k+1$; it does not force us to use it. The engines generalize: if $s=2^a+d$ with
$\max_{0\le y\le d}s_2(y)\le a-2$, then the unique maximizer of $s_2(s-t)$ is $s-t=2^a-1$ (any
$x=2^a+y$, $y\le d$, has $s_2=1+s_2(y)\le a-1$), so \emph{$\pi_{d+1}$ odd forces
$v_2(\perm(2e+J))=s-a$ exactly} (the case $d=0,1$ recovers
Theorems~\ref{thm:engineA}--\ref{thm:engineB}; the side condition is sharp---$a=3$, $d=3$ fails via the
tie $s_2(7)=3$). Consequently, for \emph{every} odd $c$ with $d=(c-1)/2$ small and $k\ge5$, the peak
$n=2^k+c$ has factor sizes $2^{k-1}+d$ and $2^{k-1}+d+1$, both engine-covered, and
\[
   D_{N_0}(2^k+c)=2k-3-s_2(c)
   \iff
   \#\{v:\pi_{d+1}(A_v),\ \pi_{d+2}(B_v)\ \text{both odd}\}\ \text{odd}.
\]
Moreover \emph{all} the parities $\pi_j\bmod2$ are computable with no permanents: by the mod-4
factorization identity the $\chi_4$-matrix decomposes over $\mathbb F_2$ as a sum of rank-one divisor
indicators $e=\sum_{q^a}\mathbf1_{q^a\mid\cdot}\mathbf1_{q^a\mid\cdot}^{\mathsf T}$, and Cauchy--Binet
plus M\"obius inversion on the partition lattice (blocks of size $\ge3$ carry even $(|B|-1)!$) give
\[
   \pi_j(e)\ \equiv\ \sum_{|T|=j}a_T^{\mathrm{row}}\,a_T^{\mathrm{col}}\pmod2,\qquad
   a_T=\!\!\sum_{\substack{\text{partitions of }T\\ \text{into singletons/pairs}}}\ \prod_B
   N(\mathrm{lcm}\,B),
\]
with $N(D)$ a $\lfloor m/D\rfloor$-type multiple count---pure divisor combinatorics (validated against
brute-force minors and against all fixed-$c{=}1$ counts). Computing the resulting \emph{adaptive-peak
table} for $k\le10$, $c\le7$: every $k$ has an odd-parity entry with $c\le3$; in particular the
$c{=}1$ failures $k=6,9,10$ are all repaired at $c=3$. The case $(k,c)=(6,3)$ is confirmed
\emph{out of sample} by an independent exact computation: the theory demands
$D_{N_0}(67)=2\cdot6-3-s_2(3)=7$, and the measured value is $7$. The open step is thus weakened from a
single-sequence parity to a window disjunction:

\begin{conjecture}[window]\label{conj:window}
There is $C_0$ (empirically $C_0=3$ suffices for all $k\le10$) such that for infinitely many $k$ some
odd $c\le C_0$ has odd achiever parity. Equivalently $D_{N_0}(2^k+c)\ge2k-3-s_2(c)\ge 2k-5$ along a
sequence of adaptive peaks, so the deficit of $N_0$ is unbounded.
\end{conjecture}

\noindent Conjecture~\ref{conj:window} is still a parity statement---the barrier is weakened, not
destroyed---but its failure would require the parity to vanish \emph{simultaneously across the whole
window} for all large $k$, a conspiracy across several independent-looking divisor correlations, and
each instance is a cheap divisor computation.

\paragraph{The collapsible family.} One further family exists on the other side: $n=2^k-1$, whose
$A$-factor has size $2^{k-1}-1$ and fires \emph{unconditionally} (the extremal term is $t=0$ with
$\pi_0=1$). The achiever set is then cut out by the single condition $\pi_1(B_v)$ odd, and---unlike
the two-condition families---its count parity \emph{collapses}: the number of odd $v$ is even, so the
constant term drops, and swapping the $v$-sum to the divisor side gives (with $r_D=s_D$ here)
\[
   C_{-1}(k)\ \equiv\ \#\{\,q^a\le n,\ q\equiv3\ (4)\ :\ \lfloor n/q^a\rfloor\equiv1,2\ (\mathrm{mod}\
   4)\,\}\ \pmod2,
\]
the parity of an alternating count of prime powers in the harmonic intervals
$\bigcup_j(n/(4j{+}3),\,n/(4j{+}1)]$. This parity \emph{provably vanishes}:

\begin{theorem}[the collapsible family never fires]\label{thm:cminus}
For $k\ge4$, the achiever count of the $n=2^k-1$ family is even; consequently
$v_2(N_0(2^k-1))$ always lifts strictly above the engine minimum.
\end{theorem}
\begin{proof}
By the collapse, the count is $\equiv\sum_{D\in P}\mathrm{oddmult}(D)\pmod 2$, where
$\mathrm{oddmult}(D)=\#\{x\ \mathrm{odd}\le n:D\mid x\}$. Swapping the order of summation
(the divisor hyperbola) gives $\sum_{x\ \mathrm{odd}\le n}\Omega_3(x)$ with
$\Omega_3(x)=\sum_{q\equiv3(4)}v_q(x)$, and by the pointwise mod-$4$ identity
$\Omega_3(x)\equiv[\,x\equiv3\ (4)\,]\pmod 2$. Hence the count is
$\equiv\#\{x\le 2^k-1:x\equiv3\ (4)\}=2^{k-2}\equiv0\pmod2$.
\end{proof}

\noindent (This matches the data: set sizes $4,6,10,18,24,44,76,126,226,410,744$ for $k=4,\dots,14$,
all even, and the measured massive lift $D_{N_0}(63)=0$.) Theorem~\ref{thm:cminus} sharpens the
obstruction picture into a proved dichotomy: \emph{linear} (single-condition) parities collapse via
the hyperbola swap to $\chi_4$-counts and vanish identically, so the deficit's unboundedness lives
\emph{exactly} in the non-collapsible correlation families of Conjecture~\ref{conj:window}.

\paragraph{The normal form: the window count is an $\mathbb F_2$ bilinear form.} The correlation
families themselves admit a drastic reduction. The key is a \emph{linearity lemma}: shifting the
column-count function by the divisor indicator $\chi_v(D)=[D\mid v]$ perturbs every matching sum only
linearly mod~$2$,
\[
   \mathrm{ms}_j(T,\,N\pm\chi_v)\ \equiv\ \mathrm{ms}_j(T,N)\ +\ \sum_{D\in T}[D\mid v]\cdot
   \mathrm{ms}_{j-1}(T\setminus D,\,N)\pmod2,
\]
proved by induction on the matching recursion: the quadratic terms $[D\mid v][E\mid v]$ arising from
the singleton product and from the pair block (via $[\mathrm{lcm}(D,E)\mid v]=[D\mid v][E\mid v]$)
occur in equal pairs and cancel mod~$2$ (verified exhaustively for $j\le4$). Hence \emph{every}
$\pi_j$-parity is affine-linear in the divisibility vector, $P(v)=b\oplus\sum_D\gamma_D[D\mid v]$ with
computable $\gamma_D=\sum_{T\ni D}\mathrm{ms}_j(T,\mathrm{row})\,\mathrm{ms}_{j-1}(T\setminus
D,\mathrm{col})$, and summing $P_1P_2$ over $v$ (hyperbola swap again) collapses the achiever count of
\emph{every} adaptive-peak family to
\begin{equation}\label{eq:normalform}
   C(k,c)\ \equiv\ b_1b_2V\ \oplus\ b_2L(\alpha)\ \oplus\ b_1L(\beta)\ \oplus\
   \alpha^{\mathsf T}M\,\beta \pmod 2,
\end{equation}
with $V=\lceil(c{+}1)/2\rceil\bmod2$, $L(\gamma)=\sum_D\gamma_D\,\mathrm{oddmult}(D)$, and kernel
matrix $M_{DE}=\mathrm{oddmult}(\mathrm{lcm}(D,E))$ over the prime powers $q^a\le n$, $q\equiv3\ (4)$.
Formula~\eqref{eq:normalform} is validated against all exactly-computed counts ($13$ families, plus
cross-engine checks). Conjecture~\ref{conj:window} is thereby a statement about a single explicit
$\mathbb F_2$ \emph{bilinear form} in two computable divisor vectors---the $v$-sum is gone---and the
per-family computation drops to one subset enumeration plus $O(|P|^2)$, eliminating the
per-$v$ phase entirely.

Three structural facts about this form (computed $k\le9$): $M$ is nearly full-rank over $\mathbb F_2$
(corank $\le5$ at $|P|=35$), so no kernel degeneracy helps; the $B$-side support $S_\beta$ is
\emph{sparse and nearly constant} ($|S_\beta|\approx7$ while $|P|$ grows), so
$\alpha^{\mathsf T}M\beta$ is a XOR of only $\sim\!7$ twisted linear sums; and the $A$-side linear
term reduces exactly (verified $k\le12$) to a \emph{digit--remainder} prime sum: with
$q=\lfloor2^{k-1}/D\rfloor$, $\rho=2^{k-1}\bmod D$,
\[
   L(\alpha)\ \equiv\ \#\{\,D\in P:\ q\ \text{odd},\ \rho<(D{-}1)/2\,\}\pmod2 .
\]
At the count level such sums are classical equidistribution of fractional parts $N/p$; at the parity
level they are Gelfond-type digit-conditions-over-primes statements---the problem family where
count-level results were achieved by Mauduit--Rivat-style exponential-sum methods. Equivalently (an
identity we verify exactly), $L(\alpha)$ is the parity of $\#\{D\in P:2^{k-1}\bmod 2D\in[D,3D/2)\}$:
where the \emph{orbit of $2$} lands in $\mathbb Z/2D$.

Moreover the $B$-side support is \emph{band-structured}: $\beta_E$ is constant on every band
$\{E:\lfloor n/E\rfloor=t\}$ with $t\ge3$ (verified $k\le10$; the exceptions at $t\le2$ are
tower-driven), so both $\alpha$ and $\beta$ are interval-indicator vectors up to bounded corrections.
The bilinear term then reduces to a sum over band-pairs of parities of counts of \emph{products}
$DE\le n$ of prime powers $\equiv3\ (4)$ in prescribed intervals. The window conjecture is thereby a
bounded combination of parities of prime- and semiprime-counting functions in explicit intervals.
This is the endpoint of the reduction, and it is two-sided: parity of prime-counting functions is
itself a famous open problem (the barrier in its most classical dress), yet
Theorem~\ref{thm:cminus} shows such combinations \emph{can} telescope to provable arithmetic-progression
counts.

\paragraph{Telescope versus residue: the verdict.} Pushing the swap calculus fully (each step
verified exactly, $k\le13$): every \emph{pure digit sum} telescopes and is provable---e.g.
$\sum_{D\in P}\lfloor 2^{k-1}/D\rfloor\bmod 2\equiv\#\{v\le2^{k-1}:\mathrm{oddpart}(v)\equiv3\
(4)\}=2^{k-2}-1\equiv1$, by the same swap-plus-pointwise-$\chi_4$ mechanism as
Theorem~\ref{thm:cminus}---and the $c{=}1$ linear term decomposes exactly as
\[
   L(\alpha)\;=\;1\ \oplus\ \underbrace{\#\{D\in P:\text{binary digits }(k{-}1,k)\text{ of }1/D
   \text{ both }1\}}_{\mathrm{Corr}(k)}\ \oplus\
   \underbrace{\#\{D\in P: D\mid 2^k{+}1,\ \lfloor2^{k-1}/D\rfloor\text{ odd}\}}_{\mathrm{Div}(k)} ,
\]
with $\mathrm{Div}$ an explicit factorization term. What resists is precise: single digits
telescope, digit \emph{pairs} do not ($\mathrm{Corr}$ is the adjacent-digit correlation of $1/D$
over $D\in P$); and for the bilinear term the $u$-swap gives (verified)
$\alpha^{\mathsf T}M\beta\equiv\#\{u\ \text{odd}\le n: \textstyle\sum_{D\mid u}\alpha_D\ \text{odd}\
\wedge\ \sum_{E\mid u}\beta_E\ \text{odd}\}$---\emph{again} a correlation of two divisor-sum
parities, the same shape as the original achiever count. The swap calculus thus strips all linear
content to provable constants and factorization terms while \textbf{fixing the correlation shape}:
correlations reproduce correlations under every available identity. The residue of
Conjecture~\ref{conj:window} is a fixed point of the reduction calculus---pairwise digit
correlations of reciprocals of primes $\equiv3\ (4)$---so progress requires genuinely second-order
input (joint digit equidistribution over primes at the parity level) or an identity outside the
swap calculus. This, and not any single computation, is the precise remaining distance to
unboundedness.

Two such outside identities exist, and they push the residue one level further (each verified
exactly). First, the digit \emph{product} is itself swappable---$q_{k-1}(D)q_k(D)$ counts pairs---so
$\mathrm{Corr}(k)\equiv\#\{(v,w):v\le2^{k-1},\,w\le2^k,\ \mathrm{oddpart}(\gcd(v,w))\equiv3\
(4)\}\pmod2$, a two-dimensional gcd-$\chi_4$ lattice count. Second, on \emph{square} domains the
transposition involution $(v,w)\mapsto(w,v)$ cancels all off-diagonal pairs mod~$2$, so with
$\kappa(a,b):=\#\{v\ \mathrm{odd}\le2^a,\ w\ \mathrm{odd}\le2^b:\gcd\equiv3\ (4)\}\bmod2$ one gets
the proven evaluation $\kappa(a,a)\equiv[a{=}2]$. A parity-quadrant recursion
($\gcd(2v,w)=\gcd(v,w)$ for odd $w$) reduces $\mathrm{Corr}(k)$ entirely to the dyadic table
$\kappa(i,j)$, whose rows are eventually periodic and whose diagonal is provable; what remains
pseudorandom is exactly the \emph{near-diagonal strips} $\kappa(a,a{+}\delta)$, $\delta$ small
(e.g.\ $\kappa(a,a{+}1)=0,0,1,1,1,0,1,1,1$ for $a=3,\dots,11$). Each identity outside the calculus
has evaluated another stratum; the surviving core is the joint distribution of $\gcd$-$\chi_4$
counts across \emph{adjacent dyadic scales}. Both provable parity strata (Theorem~\ref{thm:cminus}
and the digit-sum constant) are machine-checked
(\texttt{Paper3CMinus1.sum\_omega3\_odd\_even}, \texttt{Paper3DigitSum.sum\_omega3\_digit\_odd}),
as is the involution evaluation $\kappa(a,a)\equiv[a{=}2]$
(\texttt{Paper3Involution.kappa\_diag}).

The full recursion \emph{assembles}: chaining product-swap, quadrant splits, the involution,
dyadic strips, and the translation identity yields an exact prime-free evaluator of
$\mathrm{Corr}(k)$ in which every step is a proven identity
(\texttt{code/hyperbola\_assembly.py}; validated against brute lattice counts and against the
prime-side definition, $17/17$). Instrumenting the strata shows the involution contributions cancel
globally, leaving exactly $\mathrm{Corr}(k)=\mathrm{FP}(k)\oplus\mathrm{BD}(k)$: the full-period
stratum (the $\tau_3$/hyperbola machinery above) and the boundary stratum, whose normal form is a
sum of $\gcd$-$\chi_4$ counts over \emph{doubling windows} $(r_v,2r_v]\bmod 2v$ anchored at
$r_v=2^{j-1}\bmod 2v$. The remaining open core of Conjecture~\ref{conj:window}'s linear stratum is thus a
single explicitly-structured object: doubling-window $\gcd$ counts, a natural target for
dynamical methods.

The dynamical analysis terminates in a \emph{word normal form} (each step verified exactly). For
fixed odd $v$ the boundary bit is periodic in the scale $i$ with period $\mathrm{ord}_v(2)$, so each
$v$ carries a fixed binary word; the word is zero \emph{iff} $v$ has no prime factor $\equiv3\ (4)$
(both directions machine-checked: \texttt{Paper3CMinus1.gcd\_chi4\_empty\_of\_no\_three} and
\texttt{\_nonempty\_of\_three}); and since $\{v:\mathrm{ord}_v(2)=m\}$ consists of divisors of
$2^m-1$,
\[
   \mathrm{BD}(i,i{+}1)\;=\;\bigoplus_{v\ \mathrm{odd}\le 2^i}\ \mathrm{word}_v\big[i\bmod
   \mathrm{ord}_v(2)\big]
\]
is exactly a XOR of periodic itinerary words attached to divisors of Mersenne numbers, evaluated at
the diagonal phase (a further reflection identity, $c_v(I)=c_v(-I)$ via $w\mapsto2v-w$, constrains
the words). The periodicity itself---that any two length-$2v$ windows carry the same gcd-$\chi_4$
count, so the word is well-defined---is machine-checked (\texttt{Paper3Translation.count\_period\_eq},
via the bijection $w\mapsto w+2vk$), as is its reflection companion. No further identity among swap,
transposition, translation, and negation reduces the count; the words are the irreducible objects. The linear stratum of Conjecture~\ref{conj:window} thus rests, in final
form, on the joint irregularity of Mersenne factorizations and their divisors' $\gcd$-$\chi_4$
itinerary words: unproved, but fully mechanized, and any structural word-sum invariant would
propagate up the proven chain to the deficit itself.

\paragraph{Word theory: invariants, sub-periods, and class numbers.} The words admit their own
theory (each statement verified exactly). \emph{Word-sum invariant}: telescoping the cumulative
count along the doubling orbit gives $\sum_i\mathrm{word}_v[i]\equiv \mathrm{wt}(v)\cdot
t_3(v)\pmod2$, where $\mathrm{wt}(v)=\#\{s\in\langle2\rangle\subseteq(\mathbb Z/v)^\times:
s>v/2\}$ is the number of $1$s in the binary period of $1/v$. This invariant is now MACHINE-CHECKED (\texttt{Paper3Translation.word\_sum\_invariant}, for odd $v$ with $2^{m}\equiv1\,(v)$ and scale $\ge2$), and unconditionally in $m$ via Euler with $m=\varphi(v)$ (\texttt{word\_sum\_invariant\_totient}).
\emph{Anti-period law}: if $2^{m/2}\equiv-1\ (v)$ then $\mathrm{word}_v[i+m/2]=
\mathrm{word}_v[i]\oplus t_3(v)$ (negation identity plus even anchors; $28/28$), which with its CRT
partial-negation variant explains every observed sub-period: sub-period iff $t_3(v)$ even,
complement-word iff odd. The single-period count $t_3(v)$ has itself a clean parity, machine-checked
via the reflection involution: for odd $v$, $t_3(v)\equiv[v\equiv3\ (4)]\pmod2$
(\texttt{Paper3Negation.t3\_card\_parity}; the map $w\mapsto2v-w$ is a gcd-preserving involution of the
period whose unique fixed point $w=v$ contributes iff $v\equiv3\ (4)$). So sub-period $\iff
v\equiv1\ (4)$, complement-word $\iff v\equiv3\ (4)$. \emph{Weight evaluation}: for full-orbit primes $\mathrm{wt}(p)=(p{-}1)/2$,
so every prime $p\equiv3\ (4)$ with $2$ primitive root carries an odd-sum word; and for
\emph{half-orbit} primes $p\equiv3\ (4)$ ($2$ a quadratic residue, $\mathrm{ord}=(p{-}1)/2$), the
classical residue-distribution formula for $p\equiv3\pmod4$ --- the count of quadratic residues in
$(0,p/2)$ exceeds the expected $(p-1)/4$ by $h(-p)/2$, a standard consequence of the class number
formula \cite[\S6]{Dav} --- gives
\[
   \sum_i \mathrm{word}_p[i]\ \equiv\ \mathrm{wt}(p)\ =\ \tfrac{(p-1)/2-h(-p)}{2}\pmod 2,
\]
verified for all thirteen such $p\le359$: \textbf{the class number of $\mathbb Q(\sqrt{-p})$ governs
the word-sum invariant}. The deficit mechanism of the gcd permanent is thus, in part,
class-number-driven --- the chain \emph{deficit $\to$ achiever parity $\to$ normal form $\to$
Mersenne words $\to$ class numbers} being verified at every link and machine-checked at five.

\paragraph{The Prime-Word Theorem and $n$-gram class-number formulas.} The words themselves are
classical objects in disguise. For prime $p$: $\gcd(p,w)\equiv3\ (4)$ iff $p\mid w$ (and $p\equiv3\
(4)$), each period contains exactly one odd multiple of $p$, and window membership reduces to
(verified over all tested primes and phases)
\[
   \mathrm{word}_p[i]\;=\;\big[\,2^i\bmod p>p/2\,\big]
   \;=\;\text{the $i$-th binary digit of }1/p .
\]
\emph{The deficit boundary words of primes are the binary expansions of reciprocal primes}; composite
words decompose $\chi_4$-linearly into prime-power multiple-itineraries. Consequently the $k$-gram
statistics of a prime word are counts of the orbit of $2$ in dyadic arcs; for half-orbit primes
($p\equiv7\bmod8$, orbit $=$ quadratic residues) these are \emph{exact class-number formulas},
verified exhaustively: the digit-pair profile per period is
\[
   \#00=\tfrac{p+4h-3}{8},\quad \#01=\#10=\tfrac{p+1}{8},\quad \#11=\tfrac{p-4h-3}{8},
   \qquad h=h(-p),
\]
and the digit-triple profile (eighth arcs) involves both $h(-p)$ and $h(-8p)$:
$B_0=\tfrac{p-7+8h(-p)-2h(-8p)}{16}$, $B_1=B_2=B_4=\tfrac{p+1+2h(-8p)}{16}$,
$B_3=B_5=B_6=\tfrac{p+1-2h(-8p)}{16}$, $B_7=\tfrac{p-7-8h(-p)+2h(-8p)}{16}$ (adjacent pairs re-sum
to the quarter formulas). The negation identity underlying the anti-period law is machine-checked
(\texttt{Paper3Negation.count\_reflect}).

\begin{theorem}[the class-number ladder has depth exactly three]\label{thm:ladder}
For $p\equiv7\pmod 8$ write $S_k(p)=\#\{\mathrm{QR}\ \mathrm{in}\ (0,p/2^k)\}$. Then
\[
   S_1=\tfrac{p-1+2h(-p)}{4},\quad S_2=\tfrac{p-3+4h(-p)}{8},\quad
   S_3=\tfrac{p-7+8h(-p)-2h(-8p)}{16},
\]
while for $k\ge4$ no expression $S_k=(p+a+b\,h(-p)+c\,h(-8p)+d\,h(-16p))/2^{k+1}$ with integers
$a,b,c,d$ exists. Consequently the density, digit-pair, and digit-triple word statistics are exact
class-number formulas, and the digit-quadruple-and-finer statistics are not.
\end{theorem}
\noindent The obstruction is structural, not computational. A level-$k$ evaluation requires a real
primitive Dirichlet character of conductor $2^k$ to resolve the $2^k$-division of the arc; but
$(\mathbb Z/16)^\times\cong C_2\times C_4$ has exactly four real characters, \emph{all} of conductor
$\le 8$ ($\mathbf 1,\ \chi_{-4},\ \chi_8,\ \chi_{-8}$), and there is no real primitive character of
any conductor $2^k$ with $k\ge4$. Equivalently, the only fundamental discriminants in the family
$\chi_{-p}\cdot(\text{$2$-power})$ are $-p$ and $-8p$ ($-4p$ and $-16p$ being non-fundamental
$2^{2j}\cdot(-p)$), so $h(-p)$ and $h(-8p)$ are the only class numbers available; they suffice through
$k=3$ and nothing reaches $k=4$. Beyond the ladder the residual $|2^{k+1}S_k-p|$ is
$O(\sqrt p\log p)$ (Pólya--Vinogradov), genuine equidistribution with no closed form.
\emph{The arithmetic content of the deficit is thus completely stratified: a finite class-number core
of depth three over $\mathbb Q(\sqrt{-p})$ and $\mathbb Q(\sqrt{-2p})$, and a character-sum tail.}

A third identity type---\emph{translation}, $w\mapsto w\pm2v$ preserving $\gcd(v,w)$---splits the
strip into full periods plus boundary windows, and the full-period part telescopes (each step
verified exactly). Its foundation---that the gcd-$\chi_4$ count over \emph{any} length-$2v$ window is
the same constant $t_3(v)$, so the counting function increments by exactly $t_3(v)$ per period---is
machine-checked (\texttt{Paper3Translation.count\_window\_offset}, from the indicator's $2v$-periodicity
via the shift and reflection identities); together with the period-block count
\texttt{count\_qperiods} ($q$ periods $=q\cdot t_3$) and the strip decomposition
\texttt{count\_strip\_decomp} it forms the machine-checked telescoping core of the word-sum invariant; the range half---that the per-scale strip counts over $(2^i,2^{i+1}]$ sum to the single count over $(1,2^m]$---is also machine-checked (\texttt{count\_strip\_telescope}), as is the per-scale strip decomposition and its sum \texttt{boundary\_sum} ($\sum\mathrm{strip}=t_3\cdot\sum q+\sum\mathrm{boundary}$). These reduce the word-sum invariant, machine-checked, to $\sum_i\mathrm{word}_v[i]\equiv(\sum\mathrm{strip})\oplus t_3\cdot(\sum q)\pmod2$; the final orbit step is assembled from $\sum\mathrm{strip}\equiv0$ (\texttt{count\_full\_range\_even}: the range spans an even number of periods, via $2^m\equiv1\,(v)$) and $\sum q\equiv\mathrm{wt}(v)$ (the period-sum of $\lfloor2^{i}/2v\rfloor$ is the digit-sum of $1/v$), giving the fully machine-checked $\sum\mathrm{word}\equiv\mathrm{wt}\cdot t_3$.
Concretely: with $t=\chi_3*\varphi$ one has
$\mathrm{FP}(a)\equiv\sum_{u\le2^{a-1}}\tau_3(\mathrm{oddpart}\,u)\pmod2$, the pairing $d\mapsto n/d$
evaluates $\tau_3\bmod2$ (a square-indicator stratum for $n\equiv1$, and $\tau(n)/2$ for
$n\equiv3$), and exactly
$\sum_{n\le X,\,n\equiv3(4)}\tau(n)/2=\#\{(d,e):d\equiv1,\ e\equiv3\ (4),\ de\le X\}$---a two-AP
hyperbola count, which Dirichlet's symmetry splits into $\sqrt X$-range sums (a range-shrinking
recursion). The decisive structural consequence: \textbf{after the product-swap the residue is
prime-free}---all remaining objects are lattice counts over the integers, where the classical
hyperbola toolkit applies. The obstruction is thereby transformed from Gelfond-class digit parity
over primes into (open, but mechanical-looking) hyperbola-recursion bookkeeping; completing that
recursion, including the boundary strata, is the remaining distance to evaluating
$\mathrm{Corr}(k)$ in closed form.

\begin{thebibliography}{9}
\bibitem{Smith} H.J.S.\ Smith, \emph{On the value of a certain arithmetical
determinant}, Proc.\ London Math.\ Soc.\ (1876).
\bibitem{Valiant} L.\ Valiant, \emph{The complexity of computing the
permanent}, Theoret.\ Comput.\ Sci.\ \textbf{8} (1979).
\bibitem{Ego} G.P.\ Egorychev, \emph{The solution of van der Waerden's problem
for permanents}, Adv.\ Math.\ \textbf{42} (1981); D.I.\ Falikman (1981).
\bibitem{Ryser} H.J.\ Ryser, \emph{Combinatorial Mathematics}, Carus Mathematical
Monographs \textbf{14}, MAA, 1963.
\bibitem{Dav} H.\ Davenport, \emph{Multiplicative Number Theory}, 3rd ed.,
Graduate Texts in Mathematics \textbf{74}, Springer, 2000. (Class number formula and
the distribution of quadratic residues; P\'olya--Vinogradov.)
\bibitem{HW} G.H.\ Hardy, E.M.\ Wright, \emph{An Introduction to the Theory of
Numbers}, 6th ed., Oxford, 2008. (Mertens' theorems.)
\bibitem{MS} F.J.\ MacWilliams, N.J.A.\ Sloane, \emph{The Theory of
Error-Correcting Codes}, North-Holland, 1977.
\end{thebibliography}
\end{document}
